\documentclass[11pt,a4paper]{article}
\usepackage[utf8]{inputenc}
\usepackage{amsmath}
\usepackage{amssymb}
\usepackage{braket}
\usepackage{physics}
\usepackage{hyperref}
\usepackage{geometry}
\geometry{margin=1in}

\title{Muon Decay in Low Energy Effective Field Theory:\\
Dimension-6 Operators with Arbitrary Neutrino Flavors}
\author{}
\date{\today}

\begin{document}

\maketitle

\begin{abstract}
We present a comprehensive analysis of muon decay in the Low Energy Effective Field Theory (LEFT) framework, cataloging all dimension-6 operators relevant for $\mu^- \to e^- + \nu_\alpha + \bar{\nu}_\beta$ transitions. We include both lepton-number conserving (LNC) and lepton-number violating (LNV) operators, allowing for arbitrary neutrino flavor combinations in the final state.
\end{abstract}

\section{Introduction}

The Low Energy Effective Field Theory (LEFT) describes physics below the electroweak scale, where the $W^\pm$, $Z$, and Higgs bosons have been integrated out. The relevant degrees of freedom are leptons, light quarks, gluons, and photons. For muon decay processes, we focus on the purely leptonic sector.

The most general effective Lagrangian for muon decay can be written as:
\begin{equation}
\mathcal{L}_{\text{eff}} = \mathcal{L}_{\text{SM}} + \sum_i \frac{C_i}{\Lambda^2} \mathcal{O}_i^{(6)} + \mathcal{O}(\Lambda^{-4})
\end{equation}
where $\Lambda$ is the new physics scale, $C_i$ are Wilson coefficients, and $\mathcal{O}_i^{(6)}$ are dimension-6 operators.

\section{Lepton-Number Conserving Operators ($\Delta L = 0$)}

For the decay $\mu^- \to e^- \nu_\alpha \bar{\nu}_\beta$ with total lepton number conservation, we consider four-fermion operators with two charged leptons and two neutrinos. The flavor indices $\alpha, \beta \in \{e, \mu, \tau\}$ can take any value.

\subsection{Vector and Axial-Vector Currents}

The most relevant operators have $(V-A) \otimes (V-A)$ structure, analogous to the Standard Model:

\begin{align}
\mathcal{O}_{VV}^{\alpha\beta} &= (\bar{e}_L \gamma^\mu \mu_L)(\bar{\nu}_\alpha \gamma_\mu \nu_\beta) \\
\mathcal{O}_{AA}^{\alpha\beta} &= (\bar{e}_L \gamma^\mu \mu_L)(\bar{\nu}_\alpha \gamma_\mu \gamma_5 \nu_\beta) \\
\mathcal{O}_{VA}^{\alpha\beta} &= (\bar{e}_L \gamma^\mu \gamma_5 \mu_L)(\bar{\nu}_\alpha \gamma_\mu \nu_\beta) \\
\mathcal{O}_{AV}^{\alpha\beta} &= (\bar{e}_L \gamma^\mu \gamma_5 \mu_L)(\bar{\nu}_\alpha \gamma_\mu \gamma_5 \nu_\beta)
\end{align}

Including right-handed charged leptons:
\begin{align}
\mathcal{O}_{VV,R}^{\alpha\beta} &= (\bar{e}_R \gamma^\mu \mu_R)(\bar{\nu}_\alpha \gamma_\mu \nu_\beta) \\
\mathcal{O}_{AA,R}^{\alpha\beta} &= (\bar{e}_R \gamma^\mu \mu_R)(\bar{\nu}_\alpha \gamma_\mu \gamma_5 \nu_\beta)
\end{align}

\subsection{Scalar and Pseudoscalar Operators}

\begin{align}
\mathcal{O}_{SS}^{\alpha\beta} &= (\bar{e}_R \mu_L)(\bar{\nu}_\alpha \nu_\beta) \\
\mathcal{O}_{SP}^{\alpha\beta} &= (\bar{e}_R \mu_L)(\bar{\nu}_\alpha \gamma_5 \nu_\beta) \\
\mathcal{O}_{PS}^{\alpha\beta} &= (\bar{e}_R \gamma_5 \mu_L)(\bar{\nu}_\alpha \nu_\beta) \\
\mathcal{O}_{PP}^{\alpha\beta} &= (\bar{e}_R \gamma_5 \mu_L)(\bar{\nu}_\alpha \gamma_5 \nu_\beta)
\end{align}

\subsection{Tensor Operators}

\begin{equation}
\mathcal{O}_{T}^{\alpha\beta} = (\bar{e}_R \sigma^{\mu\nu} \mu_L)(\bar{\nu}_\alpha \sigma_{\mu\nu} \nu_\beta)
\end{equation}
where $\sigma^{\mu\nu} = \frac{i}{2}[\gamma^\mu, \gamma^\nu]$.

\subsection{Flavor Structure}

For each operator type, the indices $\alpha, \beta$ can independently take values in $\{e, \mu, \tau\}$, giving 9 flavor combinations per operator structure:
\begin{itemize}
\item $\alpha = \beta = e$: $\mu^- \to e^- \nu_e \bar{\nu}_e$
\item $\alpha = e, \beta = \mu$: $\mu^- \to e^- \nu_e \bar{\nu}_\mu$ (Standard Model)
\item $\alpha = e, \beta = \tau$: $\mu^- \to e^- \nu_e \bar{\nu}_\tau$
\item $\alpha = \mu, \beta = e$: $\mu^- \to e^- \nu_\mu \bar{\nu}_e$
\item $\alpha = \mu, \beta = \mu$: $\mu^- \to e^- \nu_\mu \bar{\nu}_\mu$
\item $\alpha = \mu, \beta = \tau$: $\mu^- \to e^- \nu_\mu \bar{\nu}_\tau$
\item $\alpha = \tau, \beta = e$: $\mu^- \to e^- \nu_\tau \bar{\nu}_e$
\item $\alpha = \tau, \beta = \mu$: $\mu^- \to e^- \nu_\tau \bar{\nu}_\mu$
\item $\alpha = \tau, \beta = \tau$: $\mu^- \to e^- \nu_\tau \bar{\nu}_\tau$
\end{itemize}

\section{Lepton-Number Violating Operators ($\Delta L = 2$)}

Lepton-number violation can occur through operators that produce two neutrinos or two antineutrinos of the same chirality. For muon decay, the most relevant LNV processes involve:

\subsection{Same-Sign Neutrino Operators}

For $\mu^- \to e^- \bar{\nu}_\alpha \bar{\nu}_\beta$ (producing two antineutrinos):

\begin{align}
\mathcal{O}_{VV}^{\alpha\beta,\text{LNV}} &= (\bar{e}_L \gamma^\mu \mu_L)(\bar{\nu}_{\alpha}^c \gamma_\mu \bar{\nu}_\beta) \\
\mathcal{O}_{AA}^{\alpha\beta,\text{LNV}} &= (\bar{e}_L \gamma^\mu \mu_L)(\bar{\nu}_{\alpha}^c \gamma_\mu \gamma_5 \bar{\nu}_\beta) \\
\mathcal{O}_{SS}^{\alpha\beta,\text{LNV}} &= (\bar{e}_R \mu_L)(\bar{\nu}_{\alpha}^c \bar{\nu}_\beta) \\
\mathcal{O}_{SP}^{\alpha\beta,\text{LNV}} &= (\bar{e}_R \mu_L)(\bar{\nu}_{\alpha}^c \gamma_5 \bar{\nu}_\beta) \\
\mathcal{O}_{T}^{\alpha\beta,\text{LNV}} &= (\bar{e}_R \sigma^{\mu\nu} \mu_L)(\bar{\nu}_{\alpha}^c \sigma_{\mu\nu} \bar{\nu}_\beta)
\end{align}

where $\nu^c = C\bar{\nu}^T$ denotes the charge conjugate, and $C$ is the charge conjugation matrix.

\subsection{Alternative LNV Channel}

For $\mu^- \to e^+ \nu_\alpha \nu_\beta$ (note the positron in final state), we would need:

\begin{align}
\mathcal{O}_{VV}^{\alpha\beta,e^+} &= (\bar{\mu}_L \gamma^\mu e_L^c)(\bar{\nu}_\alpha \gamma_\mu \nu_\beta)
\end{align}

However, this process is typically suppressed as it requires additional charge violation or contributions from other processes.

\subsection{Majorana Neutrino Mass Insertion}

If neutrinos are Majorana particles with mass term:
\begin{equation}
\mathcal{L}_m = -\frac{1}{2} m_{\alpha\beta} (\nu_\alpha^T C \nu_\beta + \text{h.c.})
\end{equation}

Then LNC operators can contribute to LNV processes through neutrino mass insertion. The effective LNV amplitude scales as $\sim \frac{m_\nu}{E_\nu}$ where $E_\nu$ is the neutrino energy.

\section{Complete Operator Basis}

The complete dimension-6 operator basis for muon decay consists of:

\subsection{Counting}

\textbf{Lepton-Number Conserving:}
\begin{itemize}
\item Vector/Axial operators: 6 structures $\times$ 9 flavor combinations = 54 operators
\item Scalar/Pseudoscalar: 4 structures $\times$ 9 flavor combinations = 36 operators
\item Tensor: 1 structure $\times$ 9 flavor combinations = 9 operators
\item \textbf{Total LNC: 99 operators}
\end{itemize}

\textbf{Lepton-Number Violating:}
\begin{itemize}
\item Same-sign neutrino operators: 5 structures $\times$ 9 flavor combinations = 45 operators
\item Note: For $\alpha = \beta$, additional factors of 1/2 may apply for identical particles
\item \textbf{Total LNV: 45 operators}
\end{itemize}

\textbf{Grand Total: 144 independent operators}

\section{Phenomenological Implications}

\subsection{Standard Model Limit}

In the Standard Model, only the operator:
\begin{equation}
\mathcal{O}_{VV}^{e\mu,\text{SM}} = (\bar{e}_L \gamma^\mu \mu_L)(\bar{\nu}_e \gamma_\mu \nu_\mu)
\end{equation}
is generated at tree level by $W$ boson exchange, with Wilson coefficient:
\begin{equation}
C_{VV}^{e\mu,\text{SM}} = \frac{4G_F}{\sqrt{2}} = \frac{4g^2}{8M_W^2}
\end{equation}

\subsection{Beyond the Standard Model}

Non-standard neutrino flavor combinations probe:
\begin{itemize}
\item Non-standard neutrino interactions (NSI)
\item Sterile neutrino mixing
\item New gauge bosons coupled to leptons differently
\item Leptoquarks (though they contribute at dim-8 for purely leptonic processes)
\end{itemize}

LNV operators constrain:
\begin{itemize}
\item Majorana neutrino masses
\item Heavy right-handed neutrino exchange
\item Lepton number violation scale
\item $R$-parity violating supersymmetry
\end{itemize}

\subsection{Current Bounds}

The muon lifetime provides stringent constraints on BSM operators. The deviation from SM prediction is typically parameterized as:
\begin{equation}
\frac{1}{\tau_\mu} = \frac{1}{\tau_\mu^{\text{SM}}}\left(1 + \sum_i \frac{C_i v^2}{\Lambda^2} + \ldots\right)
\end{equation}

Current precision measurements constrain $|C_i|/\Lambda^2 \lesssim 10^{-6}$ GeV$^{-2}$ for most operators, implying $\Lambda \gtrsim 1$ TeV for $\mathcal{O}(1)$ Wilson coefficients.

\section{Operator Matching}

These LEFT operators are matched at the electroweak scale $\mu \sim m_Z$ to the Standard Model Effective Field Theory (SMEFT). The relevant SMEFT operators at dimension-6 include:
\begin{itemize}
\item $\mathcal{O}_{ll}^{ijkl} = (\bar{L}_i \gamma^\mu L_j)(\bar{L}_k \gamma_\mu L_l)$
\item $\mathcal{O}_{Hl}^{(1),ij} = (H^\dagger i\overset{\leftrightarrow}{D}_\mu H)(\bar{L}_i \gamma^\mu L_j)$
\item $\mathcal{O}_{Hl}^{(3),ij} = (H^\dagger i\overset{\leftrightarrow}{D}_\mu^I H)(\bar{L}_i \tau^I \gamma^\mu L_j)$
\end{itemize}
where $L_i$ are lepton doublets and $H$ is the Higgs doublet.

\section{Experimental Signatures}

\subsection{Michel Spectrum}

The electron energy spectrum in muon decay (Michel spectrum) is sensitive to the Lorentz structure of the operators. Deviations from the SM prediction constrain the Wilson coefficients:
\begin{equation}
\frac{d\Gamma}{dx} \propto x^2\left[(3-2x) + \rho(4x-3) + \ldots\right]
\end{equation}
where $x = 2E_e/m_\mu$ and $\rho = 3/4$ in the SM.

\subsection{Flavor-Changing Processes}

Non-standard flavor combinations can be probed through:
\begin{itemize}
\item Correlation observables in muon decay
\item Searches for $\mu \to e \gamma$ (if accompanied by dipole operators)
\item Neutrino oscillation experiments (for neutrino NSI)
\item High-precision measurements of $\tau_\mu$ and decay distributions
\end{itemize}

\subsection{Lepton Number Violation}

LNV operators contribute to:
\begin{itemize}
\item Modified Michel spectrum (though suppressed by $m_\nu/E_\nu$)
\item Correlation with neutrinoless double beta decay ($0\nu\beta\beta$)
\item Complementary constraints from collider searches for heavy Majorana neutrinos
\end{itemize}

\section{Conclusion}

We have systematically cataloged all dimension-6 LEFT operators relevant for muon decay, allowing arbitrary neutrino flavor combinations in the final state. The complete basis consists of 99 lepton-number conserving operators and 45 lepton-number violating operators, totaling 144 independent operators.

This comprehensive framework enables systematic studies of beyond-the-Standard-Model physics in muon decay, with applications to neutrino physics, lepton flavor violation, and searches for new physics at the TeV scale and beyond.

\section*{References}

\begin{itemize}
\item Jenkins, E. E., Manohar, A. V., \& Stoffer, P. (2018). Low-energy effective field theory below the electroweak scale: Operators and matching. \textit{JHEP}, 2018(3), 016.
\item Cirigliano, V., et al. (2022). Charged Lepton Flavor Violation at the EIC. \textit{arXiv:2204.02490}.
\item Crivellin, A., et al. (2017). Lepton flavor violation in the Standard Model and beyond. \textit{arXiv:1711.10391}.
\end{itemize}

\end{document}
